\chapter{프로그램이 아니라 회로를 기술한다}\label{ch:hdl}

E부의 릴레이식 봄베를 다시 떠올려 보자. 이진 시프터의 회로도에는 ``먼저 이것을 하고 다음에 저것을 하라''가 한
줄도 없다. 있는 것은 릴레이와 선, 그리고 어느 핀이 어느 선에 닿는가뿐이다. 생성기도 핀마다 ``이 핀은 이 이름의
선에 닿는다''만 적었다(코드~\ref{lst:r-stage}). 그렇게 지은 회로는 코일에 전류가 들어오는 순간 동작하고, 한 단의
접점 26개는 한 코일 신호로 함께 넘어간다. 차례를 기다리는 부품은 없다.

FPGA의 회로를 적는 SystemVerilog 코드는 겉모습이 프로그램과 비슷하다. 변수 같은 것이 있고, \code{if}와 \code{for}가
있고, 대입문이 있다. 그러나 이 코드는 실행되는 것이 아니다. 하는 일은 프로그램보다 릴레이 회로도에 가깝다. 이 장의
요점은 두 문장이다.

\begin{keybox}
\begin{enumerate}
\item 하드웨어 기술 언어(HDL)의 코드는 ``무엇을 차례로 하라''가 아니라 ``어떤 부품이 어떻게 이어져
있는가''를 적은 것이다. 회로도를 글로 쓴 것이다.
\item 그렇게 적힌 부품들은 전원이 들어와 있는 한 모두 동시에, 항상 동작한다. 코일에 전류가 들어온 릴레이가
모두 함께 움직이는 것과 같다.
\end{enumerate}
\end{keybox}

릴레이 회로와 다른 점도 하나 있다. 닫힌 접점은 양쪽으로 전류를 보내는 길이지만, 게이트는 입력에서 출력으로만
신호를 보낸다. 이 차이가 FPGA 봄베의 모양을 정한다(\ref{sec:hdl-direction}절).

\section{코드 한 줄은 부품 하나다}

\Relay 릴레이 회로도를 읽는 법은 C부에서 보았다(\ref{sec:relay-logic}절). 두 접점을 직렬로 이으면 AND, 병렬로
이으면 OR, NC 접점은 NOT이다. 회로도는 이 접점들이 어느 선과 어느 선 사이에 놓였는지를 그린 것이다. ``지금 몇
번째 접점을 처리하는 중인가''라는 물음은 뜻이 없다. 전원에서 램프까지 닫힌 길이 있으면 램프가 켜져 있고, 접점이
바뀌면 램프도 그때 따라 바뀐다.

\Why 프로그램은 메모리에 놓인 명령의 목록이고, 처리 장치가 그 목록을 위에서부터 하나씩 읽어 수행한다.
그래서 ``지금 몇 번째 줄을 실행 중인가''라는 개념이 있다. 회로에는 명령을 읽는 주체가 없다. 릴레이 회로에 접점과
선뿐이듯이, FPGA의 회로에는 게이트와 전선뿐이다. 게이트는 입력이 바뀌면 조금 뒤에 출력이 바뀌는 물건이다.
HDL은 이 게이트와 전선의 \emph{배치도}를 글로 적는 방법이다. 코드의 한 줄은 수행할 동작이 아니라 놓일 부품이다.
도구가 HDL 코드를 읽어 실제 게이트의 연결로 바꾸는 일을 \term{합성}{synthesis}이라고 한다. 릴레이식에서 생성기가
적은 핀과 선 이름의 목록이 기판의 동선이 되는 것(\ref{ch:boardgen}장)과 같은 방향의 일이다. 적어 둔 연결이 실물의
연결이 된다.

\General 다음 한 줄을 보자.

\begin{lstlisting}[language=SV,numbers=none]
assign y = (a & b) | c;
\end{lstlisting}

프로그램으로 읽으면 ``\code{a}와 \code{b}를 AND하고, \code{c}와 OR해서, \code{y}에 넣는다''는 한 번의
계산이다. 릴레이 회로로 읽으면 접점 $a$와 $b$를 직렬로 잇고, 그 둘에 접점 $c$를 병렬로 단 길이다
(그림~\ref{fig:r-logic}의 AND와 OR를 합친 것). FPGA의 회로로 읽으면 그림~\ref{fig:g-assign}과 같은 부품 두 개와
전선이다. \code{y}는 값을 담는 상자가 아니라 OR 게이트의 출력에 붙은 전선의 이름이다. 릴레이 회로에서 램프로 가는
선에 이름을 붙인 것과 같다. \code{a}가 바뀌면 \code{y}는 언제든 따라 바뀐다. ``한 번 계산하고 끝''이 아니다.

\begin{figure}[ht]
\centering
\begin{tikzpicture}[x=1cm,y=1cm]
  % 왼쪽: 프로그램
  \node[font=\small\bfseries] at (1.6,2.3) {프로그램으로 읽으면};
  \node[blkn,minimum width=34mm,align=left,font=\ttfamily\footnotesize] (p1) at (1.6,1.4) {1: t = a \& b};
  \node[blkn,minimum width=34mm,align=left,font=\ttfamily\footnotesize] (p2) at (1.6,0.5) {2: y = t | c};
  \draw[->,cred,line width=1pt] (-0.9,1.4) -- (p1.west);
  \node[note,text=cred] at (-1.6,1.4) {지금\\여기};
  \node[note] at (1.6,-0.4) {한 줄씩, 차례로, 한 번};
  % 오른쪽: 회로
  \begin{scope}[xshift=7.2cm]
  \node[font=\small\bfseries] at (2.4,2.3) {회로로 읽으면};
  \node[and gate US,gate] (g1) at (1.6,1.2) {};
  \node[or gate US,gate] (g2) at (3.6,0.8) {};
  \draw (g1.input 1) -- ++(-0.9,0) node[left,font=\ttfamily\small] {a};
  \draw (g1.input 2) -- ++(-0.9,0) node[left,font=\ttfamily\small] {b};
  \draw (g1.output) -- ++(0.4,0) |- (g2.input 1);
  \draw (g2.input 2) -- ++(-2.9,0) node[left,font=\ttfamily\small] {c};
  \draw (g2.output) -- ++(0.8,0) node[right,font=\ttfamily\small] {y};
  \node[note] at (2.4,-0.4) {부품 두 개가 늘 거기 있다};
  \end{scope}
\end{tikzpicture}
\caption{같은 식을 읽는 두 가지 방법.}
\label{fig:g-assign}
\end{figure}

이 차이에서 곧바로 나오는 결과가 있다. 회로를 기술하는 문장은 \emph{적는 순서가 의미가 없다}. 릴레이
회로도에서 릴레이를 어느 것부터 그리든 완성된 회로가 같듯이, 다음 두 코드는 같은 회로다.

\begin{lstlisting}[language=SV,numbers=none]
assign t = a & b;          |   assign y = t | c;
assign y = t | c;          |   assign t = a & b;
\end{lstlisting}

\Ours 릴레이식의 스크램블러 유닛은 시프터 보드와 배선 카드를 앞뒤로 겹쳐 꽂고, 이웃 카드의 같은 자리
커넥터끼리 리본 케이블로 잇는다(그림~\ref{fig:r-unit}). 앞 카드의 출력 26선이 다음 카드의 입력 26선이 된다.
FPGA의 스크램블러도 같은 모양의 사슬이다. 기준 RTL의 \code{scrambler} 모듈은 그 전부가
\code{always\_comb}(조합 논리를 기술하는 블록) 하나다. 글자 하나가 오른쪽 로터, 가운데 로터, 왼쪽 로터,
반사판을 지나 다시 돌아오는 경로를 적었다.

\begin{lstlisting}[language=SV,caption={\code{rtl/src/scrambler.sv}의 본체(그리스 로터 부분 생략).},label={lst:g-scr}]
always_comb begin
  for (int c = 0; c < 26; c++) begin
    letter = c[4:0];
    x = wirings_pkg::rotor_fwd(right_id, core_r, letter);
    x = wirings_pkg::rotor_fwd(middle_id, core_m, x);
    x = wirings_pkg::rotor_fwd(left_id, core_l, x);
    ...
    x = wirings_pkg::reflect(refl_id, x);
    ...
    x = wirings_pkg::rotor_bwd(left_id, core_l, x);
    x = wirings_pkg::rotor_bwd(middle_id, core_m, x);
    x = wirings_pkg::rotor_bwd(right_id, core_r, x);
    perm[c] = x;
  end
end
\end{lstlisting}

프로그램처럼 보이지만 여기에는 ``실행''이 없다. \code{x = ...}가 일곱 번 이어진 것은 변수 \code{x}를 일곱
번 고쳐 쓰는 것이 아니라, 로터를 통과하는 회로 일곱 개를 \emph{한 줄로 이어 붙인} 것이다. 줄마다
\code{x}는 서로 다른 전선이다. 카드와 카드 사이의 리본 케이블이 저마다 다른 선인 것과 같다. 단을 나누는 방법은
다르다. 릴레이식은 로터를 밀기와 고정 배선으로 나누었고(\ref{sec:relay-shifter}절), 이 코드는 로터 하나를 한
줄로 적는다. 그러나 앞 단의 출력이 다음 단의 입력으로 이어지는 사슬이라는 점은 같다. 그러면 \code{for}는
무엇일까. 다음 절의 주제다.

\section{반복문은 시간이 아니라 공간을 쓴다}

\Relay 릴레이식 봄베에는 같은 회로를 여러 벌 놓은 곳이 많다. 이진 시프터의 한 단은 26선 각각에 전환 접점 하나씩,
같은 모양의 접점 26개다(그림~\ref{fig:r-stage}). 시프터 하나는 이런 단이 다섯이어서 전환 접점이
$5 \times 26 = 130$개이고, 가는 길과 오는 길의 시프터 한 쌍을 한 릴레이의 두 극에 실어 보드 한 장이 릴레이 130개다.
유닛 하나에 이 보드가 넉 장이다. 생성기는 이것을 반복문으로 적었다. 코드~\ref{lst:r-stage}의
\code{for i in range(N)}은 회로도를 그릴 때 돌고, 그 결과는 기판 위에 동선으로 이어진 릴레이 26개다. 반복은 그리는
쪽의 일이고, 기판에 남는 것은 복사본이다. 릴레이식 6유닛이 쓰는 릴레이 약 3,500개 가운데 3,120개가 시프터인 것도
이 복사 때문이다.

\Why 프로그램의 반복문은 같은 명령을 \emph{시간을 들여} 여러 번 수행한다. 회로에는 명령을 되풀이할
주체가 없으므로, 합성 도구는 반복문을 펼쳐서 본체를 반복 횟수만큼 \emph{복사해 나란히 놓는다}. 생성기의
반복문이 릴레이 26개를 놓은 것과 같은 일이다. 반복 횟수는 합성할 때 정해져 있어야 한다. 부품의 개수는 만들기
전에 정해져야 하기 때문이다. 기판을 만든 뒤에 단의 릴레이를 27개로 늘릴 수 없는 것과 같다.

\General 4비트 두 수의 비트별 XOR을 반복문으로 적으면 XOR 게이트 네 개가 된다
(그림~\ref{fig:g-unroll}). 네 번 ``돈'' 것이 아니라 네 개가 ``놓인'' 것이고, 네 개는 동시에 동작한다.

\begin{figure}[ht]
\centering
\begin{tikzpicture}[x=1cm,y=1cm]
  \node[anchor=west,font=\ttfamily\footnotesize,align=left,draw=cgray!50,fill=clight!60,inner sep=6pt] (code) at (-0.5,1.3)
    {always\_comb\\\ \ for (int i = 0; i < 4; i++)\\\ \ \ \ y[i] = a[i] \^{} b[i];};
  \draw[->,line width=1pt] (5.0,1.3) -- node[above,note] {합성} (6.4,1.3);
  \foreach \i/\yy in {0/3.1,1/1.9,2/0.7,3/-0.5} {
    \node[xor gate US,gate,minimum height=8mm] (x\i) at (9,\yy) {};
    \draw (x\i.input 1) -- ++(-0.7,0) node[left,font=\ttfamily\footnotesize] {a[\i]};
    \draw (x\i.input 2) -- ++(-0.7,0) node[left,font=\ttfamily\footnotesize] {b[\i]};
    \draw (x\i.output) -- ++(0.6,0) node[right,font=\ttfamily\footnotesize] {y[\i]};
  }
\end{tikzpicture}
\caption{반복문의 펼침. 반복 네 번은 게이트 네 개다.}
\label{fig:g-unroll}
\end{figure}

\Ours 릴레이 유닛 안에서 26선은 저마다 자기 접점을 지난다. 단마다 선 하나에 접점이 하나씩 있기 때문이다.
코드~\ref{lst:g-scr}의 \code{for (int c = 0; c < 26; c++)}도 같다. 로터 여섯 번과 반사판을 지나는
회로 한 줄이 26개 놓인다(그림~\ref{fig:g-scr26}). 26글자의 치환이 동시에 나온다. 빠르지만 공짜가 아니다.
릴레이식이 유닛마다 릴레이 520개로 값을 치르듯이, FPGA는 칩의 자원으로 치른다. 이 모듈 하나를 합성하면 FPGA의
기본 논리 소자인 LUT를 19,717개 쓴다. 왜 이렇게 큰지, 그리고 그래서 어떻게 했는지는 \ref{ch:resources}장과
\ref{ch:lesson}장에서 다룬다.

\begin{figure}[ht]
\centering
\begin{tikzpicture}[x=1cm,y=0.9cm]
  \foreach \row/\lab in {0/A,1/B,2/C} {
    \node[font=\ttfamily\small] at (-0.6,-\row) {\lab};
    \foreach \j/\t/\st in {0/R/blkb,1/M/blkb,2/L/blkb,3/반사/blkn,4/L/blkb,5/M/blkb,6/R/blkb} {
      \node[\st,minimum width=11mm,minimum height=6mm,font=\footnotesize] (n\row\j) at (\j*1.6+0.6,-\row) {\t};
    }
    \foreach \j/\k in {0/1,1/2,2/3,3/4,4/5,5/6} { \draw[->] (n\row\j) -- (n\row\k); }
    \draw[->] (-0.3,-\row) -- (n\row0);
    \draw[->] (n\row6) -- ++(1.1,0) node[right,font=\ttfamily\footnotesize] {perm[\row]};
  }
  \node at (5.4,-3) {$\vdots$};
  \node[font=\ttfamily\small] at (-0.6,-3.9) {Z};
  \foreach \j/\t/\st in {0/R/blkb,1/M/blkb,2/L/blkb,3/반사/blkn,4/L/blkb,5/M/blkb,6/R/blkb} {
    \node[\st,minimum width=11mm,minimum height=6mm,font=\footnotesize] (z\j) at (\j*1.6+0.6,-3.9) {\t};
  }
  \foreach \j/\k in {0/1,1/2,2/3,3/4,4/5,5/6} { \draw[->] (z\j) -- (z\k); }
  \draw[->] (-0.3,-3.9) -- (z0);
  \draw[->] (z6) -- ++(1.1,0) node[right,font=\ttfamily\footnotesize] {perm[25]};
  \draw[decorate,decoration={brace,amplitude=5pt}] (12.9,0.3) -- (12.9,-4.2) node[midway,right=6pt,note,align=left] {같은 회로\\26벌};
\end{tikzpicture}
\caption{\code{scrambler}의 반복문이 만드는 회로. R, M, L은 오른쪽, 가운데, 왼쪽 로터를 통과하는 회로다.}
\label{fig:g-scr26}
\end{figure}

\section{접점과 다른 점: 게이트에는 방향이 있다}\label{sec:hdl-direction}

\Relay 닫힌 접점은 그냥 전선이다(\ref{sec:relay-logic}절). 전류는 어느 쪽으로든 흐른다. 릴레이식 봄베는 이
성질을 쓴다. 유닛 끝의 $S(-l)$은 ``$l$칸 밀기'' 보드를 거꾸로 끼워 만들어 뺄셈 회로가 필요 없고
(\ref{sec:relay-shifter}절), 메뉴에 건 전압은 스크램블러를 앞으로도 거꾸로도 지나 퍼진다(\ref{sec:f-light}절).
대가는 샛길이다. 의도하지 않은 우회로로도 전류가 흐르므로, 시프터는 한 단의 접점을 모두 함께 움직이고 테스트
레지스터는 선마다 다이오드를 넣어 막는다(\ref{sec:t-register}절).

\Why 게이트는 다르다. 게이트의 출력은 입력이 정하고, 출력 쪽 전압이 입력을 바꾸지는 않는다. 그래서 접점 하나가
하던 일을 게이트로 하려면 게이트가 둘 든다. ``P가 켜져 있고 접점이 닫혀 있으면 Q를 켠다''와 ``Q가 켜져 있고
접점이 닫혀 있으면 P를 켠다''다. 그런데 이 둘을 그대로 이으면 P에서 Q로 간 신호가 다시 P로 돌아오는 고리가
생긴다(그림~\ref{fig:g-direction}). 접점이 닫혀 있는 동안 두 게이트는 서로를 붙들고 있어서, 처음 P를 켠 전압을
거두어도 꺼지지 않는다. 릴레이 회로에서는 전압을 거두면 전류가 사라지지만, 게이트의 고리는 의도하지 않은 기억이
된다.

\General 그림~\ref{fig:g-direction}의 왼쪽은 접점 $k$ 하나로 이은 두 선이고, 오른쪽은 같은 일을 게이트로 적은
것이다. 왼쪽에서 방향은 정할 필요가 없다. 오른쪽에서는 방향마다 회로가 따로 있다.

\begin{figure}[ht]
\centering
\begin{circuitikz}[x=1cm,y=1cm]
  \ctikzset{bipoles/length=9mm}
  % 왼쪽: 닫힌 접점 하나
  \node[font=\small\bfseries] at (1.4,2.5) {릴레이: 닫힌 접점 하나};
  \draw (0,0.7) node[left,font=\ttfamily\small] {P} to[ncs,l={\scriptsize $k$}] (2.8,0.7) node[right,font=\ttfamily\small] {Q};
  \draw[->,cblue,line width=0.8pt] (0.5,1.55) -- node[above,note,text=cblue] {P $\to$ Q} (2.3,1.55);
  \draw[->,cblue,line width=0.8pt] (2.3,-0.1) -- node[below,note,text=cblue] {Q $\to$ P} (0.5,-0.1);
  % 오른쪽: 게이트 두 개
  \begin{scope}[xshift=7cm]
  \node[font=\small\bfseries] at (2.1,2.5) {게이트: 방향마다 하나};
  \coordinate (pb) at (0,0); \coordinate (qb) at (4.2,0);
  \node[and gate US,gate] (d1) at (2.1,1.3) {};
  \node[and gate US,gate,rotate=180] (d2) at (2.1,-0.1) {};
  \draw[cred,line width=0.9pt] (d1.input 1) -- (d1.input 1 -| pb);
  \draw (d1.input 2) -- ++(-0.4,0) node[left,font=\ttfamily\scriptsize] {k};
  \draw[cred,line width=0.9pt] (d1.output) -- (d1.output -| qb);
  \draw[cred,line width=0.9pt] (d2.input 1) -- (d2.input 1 -| qb);
  \draw (d2.input 2) -- ++(0.4,0) node[right,font=\ttfamily\scriptsize] {k};
  \draw[cred,line width=0.9pt] (d2.output) -- (d2.output -| pb);
  \draw[cred,line width=0.9pt] (d1.input 1 -| pb) -- (d2.output -| pb);
  \draw[cred,line width=0.9pt] (d1.output -| qb) -- (d2.input 1 -| qb);
  \node[font=\ttfamily\small,anchor=east] at (-0.05,0.6) {P};
  \node[font=\ttfamily\small,anchor=west] at (4.25,0.6) {Q};
  \node[note,text=cred] at (2.1,-1.2) {고리: P $\to$ Q $\to$ P};
  \end{scope}
\end{circuitikz}
\caption{접점 하나와 같은 일을 하는 게이트. 접점 $k$가 닫혀 있으면 전류는 양쪽으로 흐른다. 게이트는 한쪽으로만
보내므로 방향마다 AND 게이트가 하나씩 들고, 둘을 이으면 빨간 고리가 생긴다. $k$가 1인 동안 P와 Q는 서로를
붙들어, 처음 P를 켠 전압을 거두어도 꺼지지 않는다.}
\label{fig:g-direction}
\end{figure}

\Ours 봄베의 배선망 전체가 이런 접점의 덩어리다. 그대로 게이트로 옮기면 메뉴 곳곳에 고리가 생기고, 다음 위치를
시험하려고 선을 지울 방법이 없다. 그래서 FPGA 봄베는 고리를 레지스터로 끊는다. 676가닥의 상태를 레지스터에 담고,
클럭이 한 번 뛸 때마다 켜진 선 전부가 한 걸음씩 퍼지게 한다(\ref{sec:f-rounds}절). 다음 절의 전파 코드에서
\code{if}가 두 줄씩 짝지어 있는 것은 방향마다 하나씩이기 때문이다. 레지스터와 클럭은 \ref{ch:sync}장에서, 고리를
끊는 결정 자체는 \ref{ch:rtl-core}장에서 본다.

\section{모든 것이 동시에 동작한다}\label{sec:concurrent}

\Relay 릴레이는 코일에 전류가 들어오는 그때 움직인다. 다른 릴레이의 차례를 기다리지 않는다. 시프터 한 단의 접점
26개는 한 코일 신호로 함께 넘어가고(\ref{sec:relay-shifter}절), 계수기들은 같은 \code{CK} 펄스 안에서 모두 함께
움직인다(\ref{sec:relay-counter}절). \code{STROBE}가 시험 전압을 메뉴 배선망에 걸면(\ref{sec:t-register}절) 전류는
메뉴의 모든 유닛으로 한꺼번에 퍼진다. 릴레이식에서 차례가 필요한 일, 곧 전압을 걸고, 판정하고, 지우고, 다음
위치로 가는 일은 순서기라는 장치가 따로 만든다(\ref{ch:relay-sequencer}장).

\Why 부품은 전원이 들어와 있으면 늘 동작한다. 회로의 한쪽이 일하는 동안 다른 쪽이 기다리는 일은 없다.
FPGA도 같다. 코드에 블록이 여러 개 있으면 그 블록들은 \emph{차례로} 실행되는 것이 아니라 \emph{나란히} 놓여
동시에 동작한다. 프로그램에서 여러 일을 동시에 하려면 스레드를 따로 만들어야 하지만, 회로에서는 동시가
기본이고 오히려 ``차례로''를 만들려면 따로 장치를 두어야 한다. 릴레이식의 순서기에 해당하는 그 장치가 FPGA에서는
상태 기계다(\ref{sec:fsm}절).

\General LED를 깜박이면서 동시에 버튼이 눌린 횟수를 세는 회로를 생각하자.

\begin{lstlisting}[language=SV,numbers=none]
// (1) blink: a free-running counter, the LED shows its top bit
always_ff @(posedge clk) blink_count <= blink_count + 24'd1;
assign led = blink_count[23];

// (2) count button presses
always_ff @(posedge clk) if (button_pressed) presses <= presses + 8'd1;
\end{lstlisting}

프로그램이라면 ``깜박이는 루프''에 들어가 있는 동안 버튼을 볼 수 없어서 고민해야 한다. 회로에서는 (1)과
(2)가 서로 다른 부품이므로 아무 고민이 없다. 둘은 같은 클럭에 맞춰 각자 동작한다
(그림~\ref{fig:g-concurrent}). 릴레이식의 계수기 여럿이 같은 \code{CK}에 각자 오르는 것과 같다.

\begin{figure}[ht]
\centering
\begin{tikztimingtable}[timing/slope=0.1,timing/coldist=4pt,xscale=1.7,yscale=1.25,semithick]
  \code{clk}             & 10{.5L .5H} \\
  \code{blink\_count}    & .5D{7} 1D{8} 1D{9} 1D{10} 1D{11} 1D{12} 1D{13} 1D{14} 1D{15} 1D{16} .5D{17} \\
  \code{button\_pressed} & 2.5L 1H 3L 1H 2.5L \\
  \code{presses}         & 3.5D{0} 4D{1} 2.5D{2} \\
\extracode
  \begin{pgfonlayer}{background}
    \vertlines[cgray!50,dashed]{0.5,1.5,...,9.5}
  \end{pgfonlayer}
\end{tikztimingtable}
\caption{두 블록은 같은 클럭 에지에서 각자 갱신된다. 한쪽이 다른 쪽을 기다리지 않는다.}
\label{fig:g-concurrent}
\end{figure}

\Ours 봄베에서 이 성질은 이어진 선을 찾는 방법을 정한다. A부에서 본 세 방법을 다시 놓아 보자. 실물 두 기계는
전류에 맡긴다. 전압을 걸면 이어진 선이 모두 한순간에 켜진다(\ref{sec:f-current}절). 소프트웨어(Rust)는 선을
\emph{한 가닥씩} 꺼내 그 이웃을 살핀다(\ref{sec:f-stack}절). FPGA에는 퍼지는 전류가 없고, 앞 절에서 본 대로
게이트를 고리로 이을 수도 없다. 대신 동시성을 쓴다. 회로는 메뉴의 모든 선, 양방향, 대각선 보드를 \emph{한꺼번에}
적용한다. 클럭 한 번에 통전된 영역 전체가 한 걸음씩 넓어진다(그림~\ref{fig:g-wave}, \ref{sec:f-rounds}절).

\begin{figure}[ht]
\centering
\begin{tikzpicture}[x=0.8cm,y=0.8cm,
   nd/.style={circle,draw=cgray,fill=white,minimum size=3.2mm,inner sep=0pt},
   on/.style={circle,draw=cblue,fill=cblue,minimum size=3.2mm,inner sep=0pt}]
  % 그래프 모양: 0-1,0-2,1-3,2-3,3-4,2-5,4-6
  \foreach \p/\dx/\ttl/\lit in {0/0/시작/{0},1/4.1/1클럭 뒤/{0,1,2},2/8.2/2클럭 뒤/{0,1,2,3,5},3/12.3/3클럭 뒤/{0,1,2,3,4,5}} {
    \begin{scope}[xshift=\dx cm]
      \coordinate (c0) at (0,0); \coordinate (c1) at (1,1); \coordinate (c2) at (1,-1);
      \coordinate (c3) at (2,0); \coordinate (c4) at (3,0.8); \coordinate (c5) at (2,-1.8); \coordinate (c6) at (4,0.2);
      \foreach \a/\b in {0/1,0/2,1/3,2/3,3/4,2/5,4/6} { \draw[cgray] (c\a) -- (c\b); }
      \foreach \i in {0,...,6} { \node[nd] at (c\i) {}; }
      \foreach \i in \lit { \node[on] at (c\i) {}; }
      \node[font=\small] at (2,2.0) {\ttl};
    \end{scope}
  }
\end{tikzpicture}
\caption{회로에서의 확산(개념도). 점은 선 한 가닥, 선분은 배선이다. 클럭마다 통전된 영역의 가장자리 전체가
동시에 한 걸음 나아간다. 소프트웨어는 같은 영역을 점 하나씩 채우고, 실물 기계의 전류는 한순간에 채운다.}
\label{fig:g-wave}
\end{figure}

기준 RTL에서 이 ``한꺼번에''를 만드는 것도 반복문이다. 다음 코드는 \code{bombe\_core}의 전파 논리다.
바깥 반복은 메뉴의 선, 안쪽 반복은 26가닥의 선을 돈다. 앞에서 본 대로 이 반복문은 시간을 쓰지 않는다.
메뉴의 선 16개 × 글자 26개 × 양방향에 해당하는 회로가 전부 나란히 놓인다. 안쪽의 \code{if} 두 줄이 그
양방향이다. 한 줄은 케이블 $a$에서 $b$로, 다른 한 줄은 $b$에서 $a$로 켠다(\ref{sec:hdl-direction}절).

\begin{lstlisting}[language=SV,caption={\code{rtl/src/bombe\_core.sv}의 전파 논리(일부).},label={lst:g-spread}]
always_comb begin
  live_n = live_q;
  for (int i = 0; i < MAX_EDGES; i++) begin
    cable_a = 26 * edge_a[i];
    cable_b = 26 * edge_b[i];
    if (edge_en[i]) begin
      for (int x = 0; x < 26; x++) begin
        y = {27'd0, perm_q[(i*26 + x)*5 +: 5]};
        if (live_q[cable_a + y]) live_n[cable_b + x] = 1'b1;
        if (live_q[cable_b + y]) live_n[cable_a + x] = 1'b1;
      end
    end
  end
  ...
end
\end{lstlisting}

이 회로가 FPGA가 치르는 값이다. 실물 기계에서는 전류가 공짜로 하던 일을 FPGA는 회로를 지어서 한다
(\ref{sec:f-cost}절). 그래도 세 방식은 같은 답을 낸다. 구하는 것이 순서와 무관하기 때문이다. 찾는 것은 ``전압을
건 선과 이어진 선 전체''라는 하나의 집합이다. 한순간에 찾든, 한 가닥씩 찾든, 한 겹씩 찾든 그 집합은
같다(\ref{sec:f-same}절). 이 코드는 G부에서 한 줄씩 다시 읽는다.

\section{정리}

\begin{center}
\small
\begin{tabularx}{\linewidth}{l >{\raggedright\arraybackslash}X >{\raggedright\arraybackslash}X >{\raggedright\arraybackslash}X}
\toprule
 & 프로그램 & 릴레이 회로 & HDL로 적은 회로 \\
\midrule
코드(그림)의 한 줄 & 수행할 명령 & 놓일 접점 & 놓일 부품 \\
적는 순서 & 실행 순서 & 의미 없음 & 의미 없음 (조합 논리의 경우) \\
변수 & 값을 담는 메모리 & 선의 이름 & 전선의 이름 \\
반복문 & 시간을 들여 되풀이 & 생성기가 같은 단을 여러 벌 그림 & 부품을 복사해 나란히 놓음 \\
여러 블록 & 차례로 & 동시에 & 동시에 \\
연결의 방향 & --- & 양방향 (닫힌 접점) & 입력에서 출력으로만 \\
``차례로''를 만드는 것 & 기본으로 차례다 & 순서기 & 상태 기계 \\
비용 & 실행 시간 & 릴레이의 수와 붙는 시간 & 부품의 수와 신호의 지연 \\
\bottomrule
\end{tabularx}
\end{center}

표의 ``변수는 전선의 이름''에는 예외가 있다. 값을 \emph{기억}하는 부품인 레지스터다. 릴레이식에서 자기 유지
릴레이(\ref{sec:relay-latch}절)와 계수기의 마스터--슬레이브(\ref{sec:relay-counter}절)가 하던 일이다. 다음 장에서
다룬다.
